\documentclass[11pt]{article}

%%%%%%%%%%%%%%%%%%
% Basic packages %
%%%%%%%%%%%%%%%%%%

\usepackage[T1]{fontenc}
\usepackage[utf8]{inputenc}
\usepackage[british]{babel}
\usepackage[stretch = 25, shrink = 25, tracking=true, letterspace=30]{microtype}

\usepackage{geometry}
 \geometry{
 a4paper,
 %total={150mm,257mm},
 left=25mm,
 right=25mm,
 top=15mm,
 } % Page Layout
\usepackage[dvipsnames]{xcolor} % Colour
\usepackage{amsmath} % Maths

%%%%%%%%%%%%%%%%%%%
% Tables, figures %
%%%%%%%%%%%%%%%%%%%

\usepackage{float}
\usepackage{graphicx} % Imgages, figures
\graphicspath{{./figures/}} % Set figure path

\usepackage{booktabs} % Pretty tables

%%%%%%%%%%%%%
% Citations %
%%%%%%%%%%%%%

\usepackage[
    backend=biber,
    style=apa,
    sortlocale=de_DE,
    natbib=true,
    url=false, 
    doi=true,
    eprint=false
]{biblatex}

%\addbibresource{ASSIGNMENT_1.bib} %Imports bibliography file
\usepackage{csquotes} % used for block quotes

%%%%%%%%%
% Fonts %
%%%%%%%%%

%%% Helvetica:
%\usepackage[scaled=0.95]{helvet} % Imports Helvetica, which is similar to Arial
%\renewcommand{\rmdefault}{phv} % Use Helvetica

%%% Atkinson Hyperlegible, nerdy font 
%\usepackage[sfdefault]{atkinson} %% Option 'sfdefault' if the base
%% font of the document is to be sans serif.
%\usepackage[T1]{fontenc}

%%%%%%%%%%%%%%%%%%
% Header setting %
%%%%%%%%%%%%%%%%%%

\usepackage{titlesec} % pretty headers
\titleformat{\section}
  {\normalfont\large\bfseries\color{Black}} % Formatting options
  {\thesection}{1em}{}

%%%%%%%%%%%%%%%%%%
% Title settings %
%%%%%%%%%%%%%%%%%%

\title{\textcolor{NavyBlue}{\textbf{COMP527 - CA Assignment 2 - Preprocessing}}}
\author{\textbf{Pablo Ramos - 201885602}}

\setlength\parindent{10mm}

%=========%%%%%%%%%%%=========%
%=========%% BODY %%%=========%
%=========%%%%%%%%%%%=========%

\begin{document}
% Dark Mode - Comment the bottom 2 lines to get a normal white page
%\pagecolor{black}
%\color{white} % Set the default colour to white
%
%\pagenumbering{gobble} % Turns off page numbering
%
\maketitle
{\setlength\parindent{0mm} {\color{NavyBlue} \rule{\linewidth}{0.5mm}}} % Lines with colour options
\section{About the Data}

The data has 1760 sentences from various texts. A quick read of some of them shows that it's snippets from Shakespeare's Coriolanus, a New York noir novel, food-related writing and other miscellaneous or nonsensical phrases.

\section{Loading the Data}

The data was loaded as a Pandas data frame for easier clean-up and manipulation. 

\section{Cleaning the Data}

The data will be vectorised with TF-IDF. For this to work properly, it is necessary to homogenise the text. Capitalisation needs to be removed, so words like 'sword', and 'Sword' are treated equally. Also, since this algorithm treats each sentence as a vector and each word as a component of said vector, punctuation and special characters become meaningless and need to be removed thusly.

With regex formatting the \texttt{clean\_sentence} function was defined to replace unwanted characters with whitespace. Then, a \texttt{clean\_data} function was created to apply the cleaning line by line, and label any empty lines with the word ``empty''. Both functions are applied to the data frame. To test the code, when run directly, \texttt{preprocessing.py} creates a .csv file with the outcome.

\section{TF-IDF Parameters}

Two parameters were set to further filter the vocabulary before vectorisation. \texttt{min\_df=2} discards any word appearing in fewer than two sentences, as such words are likely proper nouns, typos, or overly specific terms that would not generalise across clusters. \texttt{max\_features=4500} caps the vocabulary at the 4500 most frequent terms, keeping the resulting matrix computationally manageable and reducing noise from uncommon words.

\section{Further Steps}
Normalisation and LSA needed to be applied to the vectorised data. This was done within the actual modelling code, and will thus be explained in its documentation.

%\begin{displayquote}
%    This is some text from a book.
%\end{displayquote}

%%%%%%%%%%%%%%%%%%%%%%
% Implicit citations %
%%%%%%%%%%%%%%%%%%%%%%

%\nocite{Hunter:2007}
% \nocite{Waskom2021}


 %\printbibliography

\end{document}